\documentclass[12pt,a4paper]{article}

% ---------- Packages ----------
\usepackage[a4paper,margin=1in]{geometry}
\usepackage[T1]{fontenc}
\usepackage{lmodern}
\usepackage{xcolor}
\usepackage{listings}
\usepackage{booktabs}
\usepackage{longtable}
\usepackage{array}
\usepackage{hyperref}
\usepackage{enumitem}
\usepackage{fancyhdr}
\usepackage{titlesec}
\usepackage{tocloft}
\usepackage{microtype}

% Custom heading aliases
\newcommand{\nsection}[1]{\section{#1}}
\newcommand{\nsubsection}[1]{\subsection{#1}}

% ---------- Hyperlinks ----------
\hypersetup{
    colorlinks=true,
    linkcolor=blue!50!black,
    urlcolor=blue!60!black,
    pdftitle={OS Command Injection Notes},
    pdfauthor={Cybersecurity Notes}
}

% ---------- Colors ----------
\definecolor{primary}{HTML}{17365D}
\definecolor{secondary}{HTML}{287271}
\definecolor{codebg}{HTML}{F4F6F8}
\definecolor{codeframe}{HTML}{D5DCE5}
\definecolor{commentgreen}{HTML}{477A45}
\definecolor{keywordblue}{HTML}{2457A7}
\definecolor{stringred}{HTML}{A23B3B}

% ---------- Page Style ----------
\pagestyle{fancy}
\fancyhf{}
\fancyhead[L]{\textcolor{primary}{Web Application Security}}
\fancyhead[R]{\textcolor{secondary}{OS Command Injection}}
\fancyfoot[C]{\thepage}
\setlength{\headheight}{15pt}

% ---------- Section Formatting ----------
\titleformat{\section}
  {\Large\bfseries\color{primary}}
  {\thesection}{0.7em}{}
\titleformat{\subsection}
  {\large\bfseries\color{secondary}}
  {\thesubsection}{0.7em}{}
\titleformat{\subsubsection}
  {\normalsize\bfseries\color{primary}}
  {\thesubsubsection}{0.7em}{}

% ---------- Table of Contents ----------
\setcounter{tocdepth}{3}
\renewcommand{\contentsname}{Table of Contents}
\renewcommand{\cftsecfont}{\bfseries}
\renewcommand{\cftsecpagefont}{\bfseries}
\setlength{\cftbeforesecskip}{4pt}

% ---------- Lists ----------
\setlist[itemize]{itemsep=3pt,topsep=5pt}
\setlist[enumerate]{itemsep=3pt,topsep=5pt}

% ---------- Code Listings ----------
\lstdefinestyle{terminal}{
    backgroundcolor=\color{codebg},
    frame=single,
    rulecolor=\color{codeframe},
    basicstyle=\ttfamily\small,
    keywordstyle=\color{keywordblue}\bfseries,
    commentstyle=\color{commentgreen}\itshape,
    stringstyle=\color{stringred},
    breaklines=true,
    breakatwhitespace=false,
    keepspaces=true,
    showstringspaces=false,
    columns=fullflexible,
    tabsize=4,
    framexleftmargin=5pt,
    framexrightmargin=5pt,
    framesep=5pt,
    xleftmargin=4pt,
    xrightmargin=4pt,
    aboveskip=8pt,
    belowskip=8pt
}
\lstset{style=terminal}

% Section customization
\setcounter{secnumdepth}{0}

% ---------- Document ----------
\begin{document}

% ---------- Cover Page ----------
\begin{titlepage}
    \centering
    \vspace*{2cm}

    {\Large\textcolor{secondary}{Cybersecurity notes}\par}
    \vspace{1.1cm}

    {\Huge\bfseries\textcolor{primary}{OS Command Injection}\par}
    \vspace{0.5cm}

    {\Large Vulnerability Notes and Techniques\par}
    \vspace{1.5cm}

    \rule{0.8\textwidth}{1pt}
    \vspace{1cm}

    {\large
        A concise reference covering command execution,
        injection techniques, blind detection, and prevention.
        \par}

    \vfill

    \begin{flushleft}
        \centering
        \textbf{Prepared by}\\
        \vspace{0.3cm}
        Youssef Reda
    \end{flushleft}
\end{titlepage}

% ---------- Table of Contents ----------
\tableofcontents
\newpage

% section 1
\nsection{OS Command Injection}

\vspace{0.5cm}
\nsubsection{Definition}

OS Command Injection, also known as Shell Injection, occurs
when an application passes untrusted input to an operating-system
command unsafely, allowing an attacker to alter the intended
command execution.

\vspace{0.5cm}
\nsubsection{Impact}

\begin{itemize}
    \item \textbf{Remote Code Execution (RCE):}
          Execute commands with the privileges of the vulnerable process.
    \item \textbf{Data Exposure:}
          Read files or sensitive information accessible to the process.
    \item \textbf{Application Compromise:}
          Interfere with the application or its execution environment.
    \item \textbf{Privilege Escalation and Lateral Movement:}
          These may be possible when additional permissions or reachable
          systems permit them.
\end{itemize}

The actual impact depends on process privileges, system
configuration, and reachable resources.

\vspace{0.5cm}
\nsubsection{Common Causes}

\begin{itemize}
    \item Concatenating user-controlled input into a shell command.
    \item Unsafe use of functions such as PHP \texttt{system()},
          \texttt{exec()}, or \texttt{shell\_exec()}.
    \item Relying on weak filtering or blacklists of shell metacharacters.
\end{itemize}

\vspace{0.5cm}
\nsubsection{Vulnerable Example}

Consider a PHP application that constructs a command using
an IP address supplied by the user:

\begin{lstlisting}[language=PHP]
system("ping " . $ip);
\end{lstlisting}

If \texttt{\$ip} is not handled safely, shell syntax may allow
execution beyond the intended \texttt{ping} command.

% section 2
\newpage
\nsection{How Command Injection Works}

\vspace{0.5cm}
\nsubsection{Example Scenario}

A store availability feature accepts a product ID and store ID.
The application invokes a legacy script using the supplied values:

\begin{lstlisting}
GET /stockStatus?productID=381&storeID=29
\end{lstlisting}

The application executes:

\begin{lstlisting}
stockreport.pl 381 29
\end{lstlisting}

A vulnerable implementation may concatenate the parameters
into a command:

\begin{lstlisting}[language=PHP]
system("stockreport.pl " . $productID . " " . $storeID);
\end{lstlisting}

\vspace{0.5cm}
\nsubsection{Injection Example}

An input such as the following may alter how a Unix-like shell
parses the command:

\begin{lstlisting}
& echo aiwefwlguh &
\end{lstlisting}

The resulting command may become:

\begin{lstlisting}
stockreport.pl & echo aiwefwlguh & 29
\end{lstlisting}

The shell interprets the separators as control syntax, so the
injected \texttt{echo} can execute separately from the intended
script. Output such as \texttt{aiwefwlguh} is evidence that the
injected command ran.

\vspace{0.5cm}
\nsubsection{Why a Trailing Separator Can Matter}

Compare the following inputs:

\begin{lstlisting}
& echo aiwefwlguh
\end{lstlisting}

\begin{lstlisting}
& echo aiwefwlguh &
\end{lstlisting}

Without the final separator, the remaining original argument
may become an argument to \texttt{echo}. With it, the shell
treats the remaining text as a separate command segment.
In typical Unix-like shells, \texttt{\&} runs the preceding
command in the background. Exact behavior depends on the shell.

% section 3
\newpage
\nsection{Useful Reconnaissance Commands}

\vspace{0.3cm}

After confirming command execution in an authorized environment,
basic commands can help identify the execution context.

\vspace{0.3cm}

\begin{longtable}{>{\raggedright\arraybackslash}p{0.42\textwidth}
        >{\raggedright\arraybackslash}p{0.48\textwidth}}
    \toprule
    \textbf{Purpose}                 & \textbf{Command}       \\
    \midrule
    Current user -- Linux/Windows    & \texttt{whoami}        \\
    OS details -- Linux              & \texttt{uname -a}      \\
    OS details -- Windows            & \texttt{ver}           \\
    Network configuration -- Linux   & \texttt{ifconfig}      \\
    Network configuration -- Windows & \texttt{ipconfig /all} \\
    Network connections and ports    & \texttt{netstat -an}   \\
    Running processes -- Linux       & \texttt{ps -ef}        \\
    Running processes -- Windows     & \texttt{tasklist}      \\
    \bottomrule
\end{longtable}

% section 4
\newpage
\nsection{Blind OS Command Injection}

\vspace{0.4cm}
\nsubsection{Definition}

Blind command injection occurs when a command executes but its
output is not returned in the HTTP response. Detection therefore
depends on indirect evidence rather than direct output.

\vspace{0.4cm}
\nsubsection{Detection Methods}

\begin{itemize}
    \item \textbf{Time-based detection:}
          Cause a measurable delay and compare response times.
    \item \textbf{Output redirection:}
          Write output to a file and retrieve it later, if accessible.
    \item \textbf{Out-of-Band (OAST):}
          Trigger an external DNS or HTTP interaction that can be
          observed by the tester.
\end{itemize}

\vspace{0.4cm}
\nsubsection{Time-Based Detection}

Example Linux payload:

\begin{lstlisting}
& ping -c 10 127.0.0.1 &
\end{lstlisting}

The \texttt{-c 10} option asks \texttt{ping} to send ten packets;
it does not directly specify a ten-second delay. On many Linux
systems this takes roughly ten seconds, but actual timing varies.
A consistent response delay close to the expected duration can
indicate that the injected command executed.

\vspace{0.4cm}
\nsubsection{Redirecting Command Output}

Output redirection writes command output to a file rather than
the HTTP response.

\begin{lstlisting}
command > output.txt 2>&1
\end{lstlisting}

\begin{itemize}
    \item \texttt{>} redirects standard output.
    \item \texttt{2>\&1} sends standard error to the same destination.
\end{itemize}

Whether the file can be retrieved depends on its path, permissions,
and server configuration.
The \textbf{web root} is the directory from which the web server
serves files. Common examples include:

\begin{itemize}
    \item Apache on Linux: \texttt{/var/www/html/}
    \item IIS on Windows: \texttt{C:\textbackslash inetpub\textbackslash wwwroot\textbackslash}
\end{itemize}

These are examples only; actual locations vary. Files in the web
root may be publicly accessible depending on the configuration.

\newpage
\nsubsection{Out-of-Band Application Security Testing (OAST)}

OAST detects interactions between the target and an external
server controlled by the tester. It is useful when command output
is hidden and file-based retrieval is unavailable.

\begin{lstlisting}
Injected input
     |
     v
Vulnerable application
     |
     v
Command executes
     |
     v
Outbound DNS/HTTP request
     |
     v
Tester observes interaction

\end{lstlisting}
DNS is commonly used because it is lightweight and often permitted
by network rules. However, DNS traffic can be blocked, proxied,
or cached.
An observed interaction is evidence of an outbound lookup or
request; it does not automatically prove that every stage of a
larger exploit succeeded.

\vspace{0.5cm}
\nsubsection{Limitations and Impact}

Blind variants may hide command output, errors, and exit status.
Detection depends on:

\begin{itemize}
    \item Timing differences.
    \item Application behavior.
    \item Accessible output files.
    \item External network interactions.
\end{itemize}

Impact depends on process privileges, filesystem permissions,
outbound network restrictions, and the data or systems reachable
from the application.

% section 5
\newpage
\nsection{Shell Metacharacters and Injection Techniques}

\vspace{0.5cm}
\nsubsection{Command Separators and Operators}

Shell metacharacters are interpreted as syntax rather than
ordinary text. Common operators include:

\begin{longtable}{>{\centering\arraybackslash}p{0.18\textwidth}
        >{\raggedright\arraybackslash}p{0.72\textwidth}}
    \toprule
    \textbf{Operator} & \textbf{Typical behavior}       \\
    \midrule
    \texttt{\&}       &
    In Unix-like shells, runs the preceding command in the background;
    in Windows \texttt{cmd.exe}, separates commands.    \\
    \texttt{\&\&}     &
    Runs the second command only if the first succeeds. \\
    \texttt{||}       &
    Runs the second command only if the first fails.    \\
    \texttt{|}        &
    Pipes standard output from the first command to standard input
    of the second.                                      \\
    \texttt{;}        &
    In common Unix-like shells, separates commands that run sequentially
    regardless of the preceding command's success.      \\
    Newline           &
    Can terminate one command and start another. URL-encoded line
    feed is \texttt{\%0A}.                              \\
    \bottomrule
\end{longtable}

Operator support and behavior depend on the operating system,
shell, and application context.

\vspace{0.5cm}
\nsubsection{Command Substitution}

In Unix-like shells, command substitution runs a command and
replaces the expression with its standard output.

\begin{lstlisting}
\`command\`
$(command)
\end{lstlisting}

Example:

\begin{lstlisting}
echo $(whoami)
\end{lstlisting}

The shell runs \texttt{whoami} and passes its output to
\texttt{echo}.

\newpage
\nsubsection{Quoted Context}

Input may be placed inside single or double quotes:

\begin{lstlisting}
"User Input"
'User Input'
\end{lstlisting}

When assessing a possible injection:

\begin{itemize}
    \item Identify how the application builds the command.
    \item Determine whether the input is inside a quoted string.
    \item Check whether the input must escape that context before
          shell syntax can be interpreted.
\end{itemize}

Single and double quotes behave differently, and exact rules
vary by shell.

\vspace{0.5cm}
\nsubsection{In-Band vs. Blind Injection}

\begin{itemize}
    \item \textbf{In-band:} Command output is visible in the HTTP response.
    \item \textbf{Blind:} Output is hidden; detection uses timing,
          OAST, behavior changes, or accessible output files.
\end{itemize}

% section 6
\newpage
\nsection{Shell Expansion Concepts}

\vspace{0.5cm}
\nsubsection{Environment Variables and Parameter Expansion}

Environment variables store information used by processes
and shell scripts.

\begin{longtable}{>{\raggedright\arraybackslash}p{0.2\textwidth}
        >{\raggedright\arraybackslash}p{0.7\textwidth}}
    \toprule
    \textbf{Variable} & \textbf{Purpose}                              \\
    \midrule
    \texttt{PATH}     & Directories searched for executable programs. \\
    \texttt{HOME}     & The user's home directory.                    \\
    \texttt{USER}     & Username, when defined.                       \\
    \texttt{PWD}      & Current working directory, when defined.      \\
    \bottomrule
\end{longtable}

Parameter expansion substitutes a variable reference with its
value before command execution:

\begin{lstlisting}
echo "$HOME"
\end{lstlisting}

\vspace{0.5cm}
\nsubsection{Standard Input and Text Utilities}

Standard input (\texttt{stdin}) is a stream from which a process
reads data. Many Unix utilities can process data from standard
input, allowing commands to pass information between processes.
Text utilities commonly perform tasks such as:

\begin{itemize}
    \item Converting letter case.
    \item Reversing text.
    \item Encoding and decoding data.
    \item Searching and filtering content.
\end{itemize}

\vspace{0.5cm}
\nsubsection{PATH Security}

The \texttt{PATH} variable lists directories searched for
executable programs.
An unsafe configuration can cause an unintended executable
to run when a command is invoked by name. Applications should
use trusted execution paths and avoid relying on
attacker-influenced environment settings.

% section 7
\newpage
\nsection{Preventing OS Command Injection}

\vspace{0.5cm}
\nsubsection{Preferred Defenses}

\begin{itemize}
    \item \textbf{Avoid OS commands:}
          Use native platform APIs or libraries to perform the required task.
    \item \textbf{Avoid shell invocation:}
          If an external executable is necessary, use process APIs that
          pass the executable and each argument separately.
    \item \textbf{Validate input:}
          Apply strict allowlist validation based on the expected type
          and format.
    \item \textbf{Use least privilege:}
          Run the application with only the permissions it needs.
    \item \textbf{Restrict access:}
          Limit filesystem permissions and unnecessary outbound
          network access.
    \item \textbf{Monitor execution:}
          Log and investigate unexpected processes and outbound connections.
\end{itemize}

\vspace{0.5cm}
\nsubsection{Input Validation}

Validation should match the intended data.

\begin{itemize}
    \item \textbf{Allowlist:}
          Accept only predefined permitted values when the input
          represents a fixed choice.
    \item \textbf{Numeric validation:}
          Require a valid numeric value when a number is expected.
    \item \textbf{Alphanumeric validation:}
          Restrict to letters and numbers only when the application's
          requirements allow it.
\end{itemize}

Example of allowed operation values:

\begin{lstlisting}
ping
status
version
\end{lstlisting}

Valid numeric input:

\begin{lstlisting}
12345
\end{lstlisting}

Invalid numeric input:

\begin{lstlisting}
123; command
\end{lstlisting}

\vspace{0.5cm}
\nsubsection{Avoid Blacklist-Only Filtering}

Removing or escaping selected metacharacters is not a reliable
standalone defense. Alternative encodings, shell-specific syntax,
and missed characters may bypass incomplete filters.
Prefer avoiding shell interpretation and using safe APIs.
Validation and escaping are supplementary controls, not substitutes
for safe command construction.

% section 8
\newpage
\nsection{Quick Reference}

\vspace{0.5cm}

\begin{longtable}{>{\raggedright\arraybackslash}p{0.25\textwidth}
        >{\raggedright\arraybackslash}p{0.65\textwidth}}
    \toprule
    \textbf{Topic}    & \textbf{Remember}                            \\
    \midrule
    Command injection &
    Untrusted input changes how an OS command is parsed or executed. \\
    Impact            &
    Depends on the privileges and resources available to the
    application process.                                             \\
    Metacharacters    &
    Separators, pipes, substitutions, and quoting can alter
    shell interpretation.                                            \\
    In-band injection &
    Output is visible in the HTTP response.                          \\
    Blind injection   &
    Output is hidden; use timing, OAST, behavior changes, or
    accessible files to detect it.                                   \\
    OAST              &
    An observed external interaction is evidence of a request or
    lookup, not proof of every exploit stage.                        \\
    Blacklists        &
    Fragile when used as the main defense.                           \\
    Best defense      &
    Avoid shell commands; otherwise pass arguments separately
    and validate input.                                              \\
    Least privilege   &
    Limits the damage if command execution is compromised.           \\
    \bottomrule
\end{longtable}

\end{document}